\section{Introduction}
\label{sec:intro}

Accurate and timely weather forecasting is of paramount importance for societal resilience, disaster risk mitigation, agriculture, aviation, and economic planning. For over half a century, the cornerstone of weather prediction has been Numerical Weather Prediction (NWP), which discretizes the atmosphere into three-dimensional grids and numerically integrates the governing laws of fluid dynamics, thermodynamics, and radiative transfer~\cite{bauer2015quiet}. Operational systems such as the European Centre for Medium-Range Weather Forecasts Integrated Forecasting System (ECMWF IFS) and the National Oceanic and Atmospheric Administration Global Forecast System (NOAA GFS) provide high-fidelity medium-range predictions. Nevertheless, NWP is notoriously computationally intensive, requiring supercomputers with tens of thousands of CPU/GPU cores running for hours to generate a single forecast cycle~\cite{brenowitz2018prognostic}.

In recent years, the paradigm of weather modeling has experienced a revolutionary transition with the advent of data-driven deep learning models. Milestones such as FourCastNet~\cite{pathak2022fourcastnet}, Pangu-Weather~\cite{bi2023accurate}, GraphCast~\cite{lam2023learning}, and FuXi~\cite{chen2023fuxi} have shown that deep neural networks trained on decades of historical atmospheric reanalysis (notably ECMWF ERA5~\cite{hersbach2020era5}) can produce forecasts comparable or superior to top-tier NWP models, while accelerating inference speed by three to four orders of magnitude (generating a 10-day forecast in seconds on a single GPU).

Despite these extraordinary breakthroughs, several critical obstacles hinder the adoption of state-of-the-art AI weather models in regional and resource-constrained settings:

\begin{enumerate}
    \item \textbf{High Computational Barriers to Entry:} Flagship foundation models require massive distributed supercomputing clusters for training. For instance, GraphCast~\cite{lam2023learning} was trained on a cluster of 32 Google TPU v4 devices for several weeks on a global grid. Regional meteorological centers, academic institutions, and national agencies in developing regions lack access to such infrastructure, restricting them to off-the-shelf global inference without the ability to adapt, calibrate, or retrain models locally.
    \item \textbf{Unique Meteorological Dynamics of Southeast Asia (SEA):} Southeast Asia ($11^\circ\text{S} - 28.5^\circ\text{N}$, $92^\circ\text{E} - 142^\circ\text{E}$) represents one of the most weather-vulnerable regions on Earth, characterized by severe tropical monsoons, intense typhoon tracks originating in the Western Pacific, complex island-and-peninsula land-sea contrasts, and rugged mountainous terrain (e.g., the Annamite Range in Vietnam and the Shan Plateau in Myanmar). Tropical convective precipitation occurs predominantly at sub-grid scales, where standard global reanalysis exhibits known systematic biases.
    \item \textbf{Under-Exploitation of Satellite Remote Sensing Observations:} Most data-driven weather models rely strictly on ERA5 reanalysis fields. However, satellite observations—specifically high-frequency geostationary infrared brightness temperature (such as NASA GPM MERGIR~\cite{janowiak2001mergir})—capture real-time, independent signatures of deep cloud convective cores that are physically decoupled from model-parameterized reanalysis approximations.
\end{enumerate}

To bridge these gaps, this work presents \textbf{WeatherForecast-SEA}, an efficient, end-to-end Graph Neural Network designed specifically for regional weather forecasting over Southeast Asia at $0.25^\circ$ horizontal resolution (~28 km). WeatherForecast-SEA adopts an Encoder--Processor--Decoder topology operating on a multiscale graph with bipartite grid-to-mesh and mesh-to-grid mappings. We formulate a dual-stream architecture fusing ERA5 reanalysis fields (the ``Brain'' branch) with NASA GPM MERGIR satellite observations (the ``Eye'' branch). Crucially, the entire pipeline is engineered to execute end-to-end on accessible consumer hardware (e.g., dual Nvidia T4/P100 GPUs available on free cloud platforms such as Kaggle), demonstrating the democratized feasibility of regional AI weather forecasting.

The primary contributions of this paper are summarized as follows:
\begin{itemize}
    \item We design and implement the first dedicated regional GraphCast-style GNN tailored to the Southeast Asian domain ($159 \times 201$ regular grid, 31,959 points) predicting 14 dynamic atmospheric variables.
    \item We formulate a dual-stream data fusion pipeline combining multi-level ERA5 reanalysis with NASA GPM MERGIR infrared brightness temperature to enhance convective sensitivity.
    \item We develop systems-level engineering optimizations—including a one-time float16 memory-mapped array pre-decoding scheme and vectorized sparse message passing—reducing training overhead from 11 hours to 6 hours per run without out-of-memory (OOM) failures under strict Kaggle RAM limits.
    \item We conduct rigorous multi-start autoregressive rollout evaluations over 72-hour forecast horizons on a held-out test year (2026), providing head-to-head benchmarking against Persistence, Climatology, and a competitive U-Net baseline, supported by regional spatial breakdowns and power spectral diagnostics.
\end{itemize}
